\section{The complete certificate checker}
\label{sec:implementation}

\subsection{Finite representation and scoping}

The file \file{code/checker.py} is the complete trusted program. It uses
only Python's standard JSON module. Input is a two-element JSON array
\texttt{[goal, proof]}. Acceptance means a zero exit status and the line
\texttt{ACCEPT}; an exception, nonzero exit, or exhausted resource is not
acceptance. The \texttt{assert} statement is not used.

Zero is integer $0$. Negative integers are free numerical parameters.
Positive integers are de Bruijn bound indices, starting at $1$. A term
is an integer or a triple \texttt{[n,a,h]}; the last child binds two
variables, with the state at index $1$ and counter at index $2$.
The canonical successor is \texttt{[0,a,0]}. A formula is
\texttt{[s,t]} or \texttt{[B]}, the latter binding the inner variable
$y$ at index $1$ and witness variable $x$ at index $2$.

The initial context is the single unconditional clause
\texttt{([], [1,1])}. Every scope-extending instruction appends one
clause and introduces exactly one fresh numerical parameter. Thus for
a context of length $n$, available parameters are
$-1,\ldots,-(n-1)$, the fresh parameter is $-n$, and the automatically
selected introduction witness is $1-n$. At the root it is $0$.
Clause indices are zero based. Booleans are not accepted as integers.

Both the recursor and $D$ bind two variables in their final child, whereas
an equation binds none. One traversal therefore suffices for terms,
formulas, scope checks, substitution, shifting, and normalization.
Children are checked before a computation can erase them.

\subsection{The three proof instructions}

\begin{center}
\begin{tabular}{p{0.40\linewidth}p{0.53\linewidth}}
\toprule
Packet & Meaning\\
\midrule
\texttt{['u',j,t,p]} & Use clause $j$ at $t$; $p$ proves its antecedent
when it has one.\\[3pt]
\texttt{['i',p]} & Introduce the goal's paired binder with its fixed
environmental witness; $p$ proves the required successor instance.\\[3pt]
\texttt{['e',B,F,s,t,p,q,r]} & $p$ proves $D B$; after opening at a
fresh name, $q$ proves $s=t$, and $r$ proves $N(F(s))$; export $N(F(t))$.\\
\bottomrule
\end{tabular}
\end{center}

The placeholder of $F$ is its one extra de Bruijn variable. The equality
premise is checked before either substituted instance. The second
instance must equal the outer goal, whose full syntax and scope were
already checked. All active goals are normal forms. The payload $p$
of an unconditional Use is deliberately unused; the examples put
\texttt{null} there. This is the absence of a premise, not an accepted
unverified derivation.

\paragraph{Two substitutions, in the right order.}
To reduce the successor recursor, the checker first substitutes the
predecessor counter into the \emph{outer} of its two body slots, while
protecting the inner state slot. It then substitutes the predecessor
state into the inner slot:
\[
 \operatorname{sub}\bigl(\operatorname{sub}(h,c,1),
                         N(R(c,a;i,u.h))\bigr).
\]
The first substitution shifts the inserted counter across the still-bound
state variable; the second removes that slot again. This works when
the counter and predecessor state themselves refer to surrounding
binders. Substituting the state first without a compensating shift would
capture those variables. A separate simultaneous-substitution implementation
and a named-variable oracle check this distinction.

\subsection{A bounded certificate relation}

The companion \file{code/checker_fuel.py} takes
\texttt{[goal,proof,fuel]}. It decrements a nonnegative integer budget on
each call to the common traversal and to the inference checker, rejecting
when the budget is exhausted. Every finite successful run of the main
program has a sufficient finite budget, so the two interfaces have the
same theorems. The bounded verifier supplies a primitive-recursive
finite-certificate relation; there is no need to assume that an arbitrary
normalization completes within a uniform primitive-recursive bound.
The budget and its checking code are included in that program's separate
measurements. Concrete Python recursion and memory limits still bound
the certificates a particular process can execute.

\subsection{What was verified}

The accompanying scripts generate exact certificates and replay the
standalone modules. They include the defining computations, capture-safe
substitution under both binder kinds, full formula equality transport,
an induction proof that uses its step hypothesis, and names for nonzero
terms whose introductions contain no witness fields. Negative examples
test escaping eigenparameters, a forged equality premise, false arithmetic,
illegal clauses and indices, wrong predecessor slots, malformed syntax,
and missing induction assumptions. Runs with \texttt{python -O} retain
the same acceptance checks.

The independent normalization audit uses a named-variable implementation
and a numerical evaluator, rather than calling the checker's substitution
to produce its expected answers. A separate template-completion experiment
samples 5,000 raw sorted templates; 451 have two valid instances, and 109
of those contain extra parameters erased by normalization. All checked
instances agree with their legal canonical repairs. These finite checks
support the implementation and catch concrete binding errors. The general
soundness, completeness, and bi-interpretation claims rest on the written
proofs above.

The exact counts, hashes, subprocess results, and independently computed
metrics are recorded in \file{evidence/}. The unchanged source Lean files
are provided for attribution and inspection; they are not a newly built
Lean project for this three-instruction calculus.

